\chapter{前期量子論とAI --- 連続から離散へ、データが導く革命}

\section*{本章の学習目標}
\begin{itemize}
\item 物質が「連続」か「粒子（離散）」かという歴史的論争と、ブラウン運動による決着を知る。
\item 古典物理学の金字塔「エネルギー等分配則」の美しさと、その限界を知る。
\item 黒体放射の謎と、プランクの苦肉の策である「エネルギーの量子化（離散）」を理解する。
\item アインシュタインの光量子仮説と「固体比熱の謎」の解決による、量子論の飛躍を学ぶ。
\item バルマー系列など「実験データから法則を見出す（帰納）」プロセスを理解する。
\item ボーアの原子模型により、前期量子論がいかにして完成したかを知る。
\item 現代のAI（シンボリック回帰）が、どのようにしてデータから「物理法則の数式」を再発見しているかを知る。
\end{itemize}

\section{粒子概念の萌芽：古代ギリシャからブラウン運動まで}

「この世界は、隙間なく滑らかにつながった『連続』なものなのか、それともこれ以上分割できない『ツブツブ（離散）』の集まりなのか？」

これは、人類が自然界を観察し始めて以来、数千年にわたって抱き続けてきた根源的な問いでした。

ツブツブ派の起源は、紀元前5世紀の古代ギリシャに遡ります。哲学者デモクリトスは、「万物を細かく切り刻んでいくと、最終的に絶対に分割不可能な究極の粒に行き着くはずだ」と考え、それを\textbf{「原子（アトム：ギリシャ語で『分割できないもの』）」}と名付けました。しかし、これはあくまで哲学的な思考実験に過ぎず、その後長らく、アリストテレスらが提唱した「空間には隙間がなく、連続な物質で満たされている」という考え方が主流を占めました。

\subsection{気体分子運動論と「温度・絶対零度」の正体}

19世紀に入ると、化学の分野から「物質は種類ごとの原子からできている」という証拠が集まり始めます。そして物理学においても、ジェームズ・クラーク・マクスウェルやルートヴィヒ・ボルツマンらが\textbf{「気体分子運動論」}という革新的な理論を打ち立てました。

彼らは、目に見えない無数の「ツブツブ（分子）」が猛スピードで飛び回り、壁にビシビシとぶつかっている状態をニュートン力学を使って数学的に計算しました。すると驚くべきことに、その「分子が壁にぶつかる力の合計」が気体の「圧力」と完全に一致し、さらに次のような美しい数式が導き出されたのです。

\[\frac{1}{2} m \langle v^2 \rangle = \frac{3}{2} k_B T\]

（\(m\)は分子1個の質量、\(\langle v^2 \rangle\)は分子の速度の2乗の平均、\(k_B\)はボルツマン定数、\(T\)は絶対温度）

この数式は、物理学における奇跡的な架け橋です。左辺は「分子が飛び回るスピード（ミクロな運動エネルギー）」を表し、右辺は「私たちが温度計で測る温度（マクロな熱）」を表しています。つまり、\textbf{「温度の正体とは、目に見えない分子が暴れ回るスピードそのものである」}ことが見事に証明されたのです。熱いお湯は分子が猛烈に飛び回っており、冷たい水は分子が大人しく動いている状態だということです。

さらに、この数式は「温度の究極の下限」の存在も明らかにします。古典物理の見方では、分子の動くスピード（左辺）がゼロへ近づくと、右辺の温度\(T\)もゼロへ近づき、それ以上温度は下がらなくなります。これが\textbf{「絶対零度（マイナス273.15度）」です。物理学では、この熱運動が最小になる極限をゼロの基準とした「絶対温度（ケルビン：K）」}という単位を使います。ただし、これはあくまで古典物理の直感的な説明です。後に量子力学が明らかにするように、絶対零度でも粒子が完全に静止するわけではなく、「零点振動」と呼ばれる量子ゆらぎが残ります。

\subsection{古典物理学の金字塔「エネルギー等分配則」}

ボルツマンらはさらに、温度がもたらす熱エネルギーは、分子の「上下・左右・前後への移動」や「回転」といったあらゆる動きの方向（自由度）に対して、\textbf{「\(\frac{1}{2} k_B T\)ずつ平等に分け与えられる」という法則を導き出しました。これを「エネルギー等分配則」}と呼びます。

この法則によれば、どんな物質であっても、温度を上げればその物質の持つ自由度に応じて「決まった量のエネルギーを平等に吸収する」はずです。この数式は、物質がどれくらい温まりにくいかを示す「比熱」という性質を数学的に説明し、古典物理学の美しい金字塔とされていました。

ただし、ここには後から見れば決定的な前提が隠れていました。等分配則は、各自由度がエネルギーを連続的に、いくらでも細かく受け取れることを暗黙に仮定しています。もしエネルギーの受け渡しに「最小単位」があるなら、十分なお小遣いがない自由度は一円も受け取れない、という事態が起こります。この隠れた前提こそが、黒体放射や固体比熱で古典物理を追い詰めることになります。

\subsection{物理学界の猛反発とアインシュタインの決着}

しかし、第3章でも触れた通り、当時の物理学の主流派（エルンスト・マッハら）はこのツブツブの理論に猛反発しました。「圧力や温度は、温度計や圧力計で直接測れる『連続的で確実な現実』である。わざわざ目に見えない『原子』などという架空のツブツブを仮定しなくても、熱力学の美しい方程式だけで世界は完璧に説明できる。原子は単なる便利な計算上のフィクション（作り話）だ」と激しく攻撃したのです。

この長きにわたる「連続か、ツブツブか」の論争に決定的な終止符を打ったのは、1905年のアルバート・アインシュタインでした（彼が特殊相対性理論などを発表した「奇跡の年」の業績の一つです）。

水に浮かべた微小な花粉などが、不規則にピクピクと動く「ブラウン運動」と呼ばれる現象があります。アインシュタインは、この不規則な動きが\textbf{「目に見えない無数の水分子（ツブツブ）が、四方八方から花粉にランダムに激突している証拠である」}とし、次のような数式を導き出しました。

\[\langle x^2 \rangle = \frac{RT}{3\pi \eta r N_A} t\]

一見複雑に見えますが、この数式が持つ意味は衝撃的でした。左辺の「花粉が動いた距離」や、右辺の「温度」「液体の粘り気」「時間」などは、すべて人間が普通の実験室で直接測定できるマクロな数値です。そして唯一、\(N_A\)（アボガドロ定数：物質1モルの中に含まれる分子の数）だけが、ミクロな世界の未知の数値でした。

つまり、この数式を使えば、顕微鏡で見える花粉の動きをストップウォッチで測るだけで、「目に見えない分子の数や大きさ」を逆算して正確に弾き出すことができるのです。

数年後、フランスの物理学者ジャン・ペランがこの数式通りに精密な実験を行い、分子の数を計算でピタリと証明しました。マクロな実験データが、ミクロな未知の数値をあぶり出したのです。この決定的な証拠により、物理学界はついに「物質の正体はツブツブ（原子・分子）である」と認めざるを得なくなりました。

\section{完璧な物理学の綻び：黒体放射と「連続」の限界}

物質がツブツブであることが証明された一方、同じ19世紀の終わり頃、多くの物理学者は「光やエネルギー」に関する物理学はすでに完成したと考えていました。マクスウェルの電磁気学が光を「連続した波」であると証明し、熱力学がそのエネルギーの振る舞いを解き明かしていたからです。「あとは少数の細かい謎を片付けるだけで、新しい大発見はもうないだろう」とさえ言われていました。

しかし、その「細かい謎」の一つが、古典物理学の巨大な城を根底から崩壊させることになります。それが\textbf{「黒体放射（こくたいほうしゃ）」}の問題でした。

\subsection{溶鉱炉の光をめぐる2つの数式：ヴィーンの法則とレイリー・ジーンズの法則}

鉄などの物質を高温の溶鉱炉で熱すると、最初は赤く光り、温度が上がるにつれて黄色、そして青白く輝くようになります。物理学者たちは、この「温度と光の色（波長・振動数）の関係」を、完璧だと信じていた熱力学や電磁気学を使って数式化しようとしました。

まず1896年、ドイツのヴィルヘルム・ヴィーンが、熱力学のアプローチから次のような数式（ヴィーンの放射法則）を導き出しました。

\[u(\nu, T) = a \nu^3 e^{-b\frac{\nu}{T}}\]

（\(a, b\)は定数、\(\nu\)は光の振動数、\(T\)は絶対温度）

この数式は、振動数\(\nu\)が大きい（波長が短い）青色や紫外線の領域では、実験データと見事に一致しました。しかし、振動数が小さい（波長が長い）赤外線の領域になると、現実のデータからどうしてもズレてしまうという弱点がありました。

そこで、ヴィーンとは別のアプローチとして、イギリスの物理学者レイリー卿（1900年）とジェームズ・ジーンズ（1905年）が、空間を満たす「連続な電磁波」の波の数を数え上げるという電磁気学的な正攻法で、次のような数式（レイリー・ジーンズの法則）を導き出しました。この計算の根底にも、先ほど登場した「エネルギー等分配則（あらゆる振動のモードに平等に\(k_B T\)のエネルギーが分配される）」という大原則が使われていました。

\[u(\nu, T) = \frac{8\pi \nu^2}{c^3} k_B T\]

（\(c\)は光の速さ、\(k_B\)はボルツマン定数）

この数式は、ヴィーンの公式が苦手だった「振動数\(\nu\)が小さい領域」では実験データと完璧に一致しました。ところが、この数式にはとんでもない欠陥がありました。

光は波長が短くなる（赤\(\rightarrow\)青\(\rightarrow\)紫外線）ほど、振動数\(\nu\)が大きくなります。この数式によると、エネルギーは「\(\nu\)の2乗（\(\nu^2\)）」に比例するため、振動数が大きい紫外線やX線の領域になればなるほど、計算上のエネルギーが「無限大」に向かって激しく発散してしまうのです。

これを\textbf{「紫外破綻（Ultraviolet catastrophe）」}と呼びます。もしこの数式が正しいなら、私たちがストーブのスイッチを入れた瞬間に、そこから無限の致死的な紫外線が放出されて世界が焼き尽くされてしまいますが、現実にはそんなことは起きません。

高振動数でうまくいくが低振動数で失敗する「ヴィーンの公式」。低振動数でうまくいくが、高振動数で世界を焼き尽くしてしまう「レイリー・ジーンズの公式」。当時の完璧なはずの古典物理学をいくらこねくり回しても、全帯域で実験データに一致する「一つの完璧な数式」を作ることがどうしてもできなかったのです。物理学はここにきて、根本的で致命的な限界に直面しました。

\section{プランクの苦悩と統計力学：確率への「絶望的な屈服」}

1900年10月、ドイツの物理学者マックス・プランクは、この無限大に発散してしまう問題を解決するために、ある奇妙な数式を「数合わせ」で考案しました。

彼は、高周波でうまくいく「ヴィーンの公式」と、低周波でうまくいく「レイリー・ジーンズの公式」の分母の形を見比べ、ほんのわずかな数学的な操作（「-1」を付け足すだけ）を行うことで、2つの極端な数式を滑らかに繋ぎ合わせる（内挿する）ことに成功したのです。この新しい数式は、全帯域の溶鉱炉の光の観測データと高精度に一致しました。

しかし、プランクにとって本当の地獄の苦しみはここからでした。彼は自分が作ったその「数合わせの式」が、物理学的に一体何を意味しているのか（なぜその数式になるのか）を、背後にある根本的なメカニズムから証明しなければならなかったのです。

\subsection{絶対的真理への固執とボルツマンへの嫌悪}

プランクは元々、クラウジウスが提唱した「熱力学第2法則（エントロピー増大則）」を、宇宙における絶対に例外が許されない「厳密で神聖な真理」として深く信奉していました。

そのため彼は、第3章で触れたルートヴィヒ・ボルツマンの「統計力学」のアプローチを激しく毛嫌いしていました。ボルツマンは、「エントロピーが増えるのは、単に原子が散らばるパターンの確率が高いから（たまたまそうなるだけ）だ」と主張していたからです。プランクにとって、大自然の絶対的な法則が「確率」や「サイコロの目」のような曖昧なものに支配されているなどという考えは、到底受け入れがたいものでした。

\subsection{絶望的な決断と「数える」ことの壁}

しかし、黒体放射の数式を厳密な古典物理学の枠組みだけで証明しようといくら計算しても、絶対にうまくいきませんでした。数週間、昼夜を問わず計算に没頭し、肉体的にも精神的にも限界まで追い詰められたプランクは、ついに自らの信念を曲げる決断を下します。

「それは絶望的な行動（Act of desperation）だった。私はそれまでの物理学の信念をすべて犠牲にする覚悟を決めた」

プランクは後年、このように振り返っています。彼は、自分がそれほどまでに嫌悪していたボルツマンの確率的なアプローチ（統計力学）にすがるしかなかったのです。

プランクは、光の集まりの全体的な性質（エントロピー）を計算するために、ボルツマンの究極の数式\textbf{「エントロピー\(S = k_B \log W\)（\(W\)は状態の数・確率）」}を使うことにしました。

ところが、ここで最大の壁にぶつかります。この数式を使って計算するためには、空間にあるエネルギーがどのように分配されるか、その起こりうるパターンの数（場合の数：\(W\)）を「1個、2個、3個\ldots{}」と指折り数えなければなりません。

しかし、当時の物理学の絶対的な常識では、空間を伝わる光のエネルギーは水のように「滑らかで連続的」なもの（アナログ）でした。コップの中の水を「1個、2個」と数えられないように、連続なものは無限に細かく分割できるため、その状態の数を数式上で「数え上げる」ことは原理的に不可能だったのです。

\section{プランクの苦肉の策：「エネルギーの粒」の誕生とデジタルの衝撃}

状態の数\(W\)を「数える」ために、プランクは物理学の常識を覆す、ある暴挙に出ます。

\textbf{「光のエネルギーは滑らかな連続ではなく、ある決まった最小単位の『ツブツブ（パケット）』の集まりである」}と強引に仮定したのです。エネルギーが硬貨やサイコロのように離散的（デジタル）なものであれば、それを「1個、2個」と数え上げることが可能になります。

「自然は飛躍しない（Natura non facit saltus）」というラテン語の古い格言があるほど、物理学の世界ではすべてが滑らかで連続であることが大前提でした。プランクのこの仮定は、その大前提への明らかな反逆でした。

プランクはこのエネルギーの最小単位の粒を\textbf{「量子（Quantum）」}と名付けました。そのエネルギーの大きさ（\(E\)）は、光の振動数（\(\nu\)）に比例し、次のような極めてシンプルな数式で表されます。

\[E = h\nu\]

ここで登場する\(h\)こそが、ミクロの世界を支配し、のちに量子力学の象徴となる絶対的な定数\textbf{「プランク定数」}です。この仮説によれば、光のエネルギーは\(1h\nu, 2h\nu, 3h\nu \dots\)というように階段状に「飛び飛びの値」しかとれず、\(1.5h\nu\)のような中途半端なエネルギーは宇宙に存在しないことになります。

\subsection{なぜ日常で「ツブツブ」に気づかないのか？}

なぜニュートンやマクスウェルは、エネルギーがツブツブであることに気づかなかったのでしょうか？

それは、プランク定数\(h\)が「\(6.626 \times 10^{-34}\) J・s」という、私たちの日常感覚からすれば絶望的なまでに極小の数字だからです。

テレビやスマートフォンの画面を遠くから見ると、滑らかで連続した美しい写真に見えます。しかし、顕微鏡レベルまで限界までズームして見ると、実はそれが小さな四角い「ピクセル（画素）」の集まりであることがわかります。プランク定数\(h\)は、宇宙を構成するエネルギーの「究極のピクセルサイズ」です。それが極端に小さすぎるため、マクロな視点で見ている私たちには、エネルギーが滑らかで連続な「波」にしか見えていなかったのです。

ここで注意したいのは、プランクが最初から「光は小さな粒で飛んでいる」と断言したわけではない、という点です。彼がまず行ったのは、黒体放射のエネルギーを数え上げるために、エネルギーの受け渡しを\(h\nu\)単位に区切ることでした。この控えめな仮定が、後にアインシュタインによって「光そのものも粒のように振る舞う」という大胆な光量子仮説へ発展します。

\subsection{「エネルギーの硬貨」が紫外破綻を救う理由}

エネルギーを「ツブツブ」として数え、ボルツマンの確率の計算をやり直した結果、1900年12月14日のドイツ物理学会で、プランクはついに理論的な証明をまとめた数式を発表しました（この日が、\textbf{「量子力学の誕生日」とされています）。 これが有名な「プランクの輻射公式（黒体放射の公式）」}です。

\[u(\nu, T) = \frac{8\pi h \nu^3}{c^3} \frac{1}{e^{h\nu/k_B T} - 1}\]

この数式は、なぜレイリー・ジーンズの「紫外破綻（高周波でのエネルギー無限大）」を見事に回避できるのでしょうか？ それは分母にある\textbf{「自然対数の底（\(e\)）の指数関数」}が魔法のように効いているからです。物理的な直感（エネルギーの硬貨のアナロジー）で説明してみましょう。

物質が持っている熱エネルギー（温度によるお小遣い）の総枠は、\(k_B T\)という量で決まっています。

一方、光を放つためには、プランクの法則に従って「\(E = h\nu\)という硬貨」を作らなければなりません。

振動数\(\nu\)が小さい「赤外線」は、1円玉のような安い硬貨です。温度のお小遣い（\(k_B T\)）があれば、この1円玉を大量にバラまくことができます（ここで低周波のレイリー・ジーンズの法則と一致します。つまり、低周波では「エネルギー等分配則」がちゃんと成り立っているように見えるのです）。

ところが、振動数\(\nu\)が極めて大きい「紫外線やX線」は、100万円札のような超高額紙幣になります。温度のお小遣いが足りない場合、「0.5枚だけ発行する」ことは量子力学のルール上許されません。払えないなら、発行枚数は「ゼロ」になります。

つまり、紫外線領域になると、エネルギーの粒（\(h\nu\)）を一つ作るコストがあまりにも高額になりすぎるため、物質はそもそも紫外線をほとんど放出できなくなり、「等分配則」が完全に破綻してしまうのです。数式上では、分母の\(e^{h\nu/k_B T}\)が猛烈な勢いで無限大に向かって巨大化し、式全体を強制的に「ゼロ」へと押し潰します。これにより、ストーブから致死的な紫外線が無限に放出されるという理論上の破綻は防がれたのです。

\subsection{プランクの後悔と消せない「量子」}

見事な証明を完成させたプランクでしたが、彼は生真面目で保守的な物理学者でした。彼自身の本心では、「計算の途中で一時的に粒（ピクセル）だと仮定しても、最後にピクセルの大きさ\(h\)を極限までゼロに近づければ、いつもの滑らかで連続な古典物理学に戻るだろう（数学的なトリックに過ぎない）」と考えていました。

しかし、計算式が実験データと一致するためには、どうしても粒の大きさ\(h\)をゼロにすることはできず、\textbf{「自然界のエネルギーは本当にツブツブのままでなければならない」}という決定的な事実が残ってしまいました。

プランクはこの後、なんと10年以上もの歳月を費やして、自分の立てた「量子仮説」をなんとか撤回し、古典的な連続の物理学に戻せないかと格闘し続けました。しかし、彼がパンドラの箱を開けて世に放ってしまった「量子（Quantum）」という概念は、もはや後戻りできない物理学の大革命を引き起こすことになります。

\section{アインシュタインの飛躍：光量子と「固体比熱」の謎}

プランクが「苦肉の策」だと考えていた量子のアイデアを、宇宙の絶対的な真実として真っ向から受け入れ、物理学の別の巨大な謎を次々と解き明かしたのが、アルバート・アインシュタインでした。

\subsection{光電効果の謎と「波」の理論の敗北}

当時、「光電効果」と呼ばれる不思議な現象が知られていました。金属に光を当てると、金属の中の電子が弾き飛ばされて外に飛び出してくる現象です（現在の太陽光発電やデジカメのセンサーの原理です）。

しかし、この現象を「光は連続した波である」という当時の常識で説明しようとすると、決定的な矛盾がいくつも生じました。

矛盾1：明るくしても意味がない。光が水波のような「連続した波」であるなら、光を強く（明るく）すれば波のエネルギーは大きくなるため、どんな色の光であっても、強烈な光を長時間当て続ければエネルギーが蓄積していつか電子が飛び出すはずです。しかし現実には、赤い光（振動数が小さい光）をいくら強烈に、何時間当て続けても電子はピクリとも動きませんでした。

矛盾2：色（振動数）だけが鍵を握る。逆に、極めて弱い（暗い）光であっても、光を青や紫外線（振動数が大きい光）に変えた途端、光を当てた「瞬間に」時間遅れなく電子がポンポンと飛び出したのです。さらに、飛び出してくる電子の勢い（運動エネルギー）は、光の明るさではなく「光の色（振動数）」だけで決まっていました。

\subsection{光量子仮説：「光は空間を飛ぶ粒である」}

アインシュタインは、プランクの「エネルギーの粒」というアイデアを極限まで押し広げることで、この謎を鮮やかに解決しました。

プランクはあくまで「物質が光を吸収・放出する瞬間だけ、エネルギーのやり取りがツブツブになる」と遠慮がちに考えていました。しかしアインシュタインは、\textbf{「光の波そのものが、実は空間を飛ぶ粒子（光量子：フォトン）の集まりである」}と真っ向から主張したのです。

アインシュタインは、光電効果を「1個の光の粒（フォトン）が、金属の中の1個の電子に激突するビリヤード」のようなものだと考えました。フォトンの1粒のエネルギーは、プランクの公式通り\(E = h\nu\)で決まります。

この見方では、光の「明るさ」と「色」の役割がはっきり分かれます。明るい光とはフォトンの数が多い光であり、青い光や紫外線とは一粒あたりのエネルギー\(h\nu\)が大きい光です。だから、どれほど強い赤い光を当てても、フォトン一粒のエネルギーが足りなければ電子は飛び出せません。一方、十分に高い振動数の光なら、弱い光でも電子を飛び出させることができます。

赤い光の粒子は振動数\(\nu\)が小さいため、1粒のエネルギー（\(h\nu\)）が金属から電子を引き剥がすための最低限のバリア（仕事関数：\(W\)）を越えられません。1粒の威力が弱すぎるため、赤い光の粒をいくら大量に（明るく）マシンガンのように浴びせても、電子を弾き出すことは絶対にできないのです。

一方、青や紫外線の粒子は振動数\(\nu\)が大きいため、1粒のエネルギーがバリア（\(W\)）を余裕で超えます。そのため、光が暗く（粒の数が少なく）ても、激突した瞬間に電子を弾き飛ばすことができるのです。

アインシュタインは、飛び出す電子の運動エネルギー（\(K\)）を次のような極めてシンプルな数式で表しました（光電効果の方程式）。

\[K = h\nu - W\]

（飛び出す電子のエネルギー\(K\)は、ぶつかった光の粒のエネルギー\(h\nu\)から、脱出するための参加費\(W\)を引いた残りのお釣りである）

この数式は、のちにアメリカの物理学者ロバート・ミリカンの精密な実験によって完璧に正しいことが証明されました（皮肉なことにミリカンは光量子仮説を嫌悪しており、アインシュタインが間違っていることを証明しようと10年も実験を重ねた結果、逆にその数式の完璧さを実証してしまったのです）。

第4章のマクスウェル方程式によって「光は波である」と完全に証明されたはずなのに、アインシュタインは「光は粒子でもある」と証明してしまったのです。この\textbf{「波と粒子の二重性」}という大いなる矛盾こそが、量子力学という深淵への入り口でした。

\subsection{デュロン＝プティの法則の破綻と「固体の比熱の謎」}

さらにアインシュタインは、光だけでなく\textbf{「物質そのものの熱エネルギー」に関する古典物理学の矛盾にもメスを入れました。それが「固体の比熱（温まりにくさ）の謎」}です。

8.1節で見た古典力学の金字塔「エネルギー等分配則」によれば、固体を構成する原子の振動に対しても、熱エネルギーは温度に比例して平等に分配されるはずです。この法則に従えば、どんな固体であっても「1度温度を上げるのに必要なエネルギー（モル比熱）」は、温度に関わらず常に一定の値（約\(3R\)）になるという\textbf{「デュロン＝プティの法則」}が導かれます。

実際、常温付近では多くの金属がこの法則に綺麗に従うため、古典物理学の勝利と思われていました。

しかし、ダイヤモンドのような極端に固い物質や、金属を極低温まで冷やした場合、\textbf{「温度が下がるにつれて比熱が急激に下がり、ゼロに近づいていく」}という、エネルギー等分配則を完全に裏切る不可解な現象が観測されていたのです。

\subsection{アインシュタインモデル：熱振動も「ツブツブ」である}

アインシュタインは、この比熱の謎も「量子（ツブツブ）」で解決できると見抜きました。

固体の原子は、バネで繋がったように特定の振動数（\(\nu\)）で揺れています。アインシュタインは、この原子の振動エネルギーも滑らかな連続ではなく、プランクの法則に完全に従って\textbf{「\(E = h\nu\)というツブツブの硬貨（フォノン：音響量子）」}でしか受け渡しできないと仮定したのです（アインシュタインの固体モデル）。

高温（お小遣い\(k_B T\)が豊富）のときは、高い硬貨（\(h\nu\)）をいくらでも払えるため、エネルギーは古典力学の等分配則通りに平等に分配され、デュロン＝プティの法則が成り立ちます。

しかし、極低温（お小遣い\(k_B T\)が極端に少ない）になると事態は急変します。原子を一段階激しく揺らすための「最初の1枚の硬貨（\(h\nu\)）」が高額すぎて、少ないお小遣いでは全く手が届かなくなってしまうのです。エネルギーの硬貨を一枚も受け取れないということは、熱を加えても原子が振動のエネルギーを吸収できない（＝比熱が急激に下がる）ことを意味します。

ダイヤモンドの比熱が常温でも極端に低いのは、原子同士の結合が強すぎて振動数\(\nu\)が大きく、硬貨（\(h\nu\)）が100万円札のように高額すぎるため、常温のお小遣すら受け取れないからなのです。

光だけでなく、物質の熱振動（力学的なエネルギー）までもが「量子化（デジタル化）」されていた。アインシュタインによるこの劇的な謎解きにより、プランクの「苦肉の策」は、単なる数合わせのトリックではなく、宇宙のあらゆるエネルギーを支配する絶対的な真理として物理学界に受け入れられていくことになったのです。

\section{スペクトルの暗号解読：データ駆動と「系列」の発見}

光の波と粒子の問題と並行して、もう一つ物理学者を悩ませていた大きな謎がありました。それが「原子のスペクトル」です。

水素などのガスをガラス管に入れて電圧をかけると、光を放ちます。その光をプリズムで分光すると、虹のような連続したグラデーションにはならず、\textbf{「特定の波長（色）の、飛び飛びの線（輝線）」}しか現れませんでした。まるで原子が発するバーコードのようなこの「飛び飛びの線」がなぜ生じるのか、誰にも分かりませんでした。

\subsection{バルマーの帰納的発見と「謎の数式」}

1885年、スイスの数学教師ヨハン・バルマーは、物理学の理論を一切使わず、ひたすらこの水素のスペクトルの波長データ（数字の羅列）を眺めていました。そして彼は、パズルのように数字をこねくり回すうちに、それらの波長が次のような「美しい分数の数式」に完璧に当てはまることを発見しました（バルマー系列）。

\[\lambda = 364.5 \times \frac{n^2}{n^2 - 2^2} \quad (n = 3, 4, 5 \dots)\]

これは純粋な\textbf{「データ駆動（帰納的）のパターンの発見」}でした。なぜ分母に\(2^2\)（4）という数字が現れるのか、物理的な「理由（演繹的な理論）」は全く分からないまま、ただ実験データだけが「美しい数学的な法則」の存在を指し示していたのです。

\subsection{未知のデータを「予言」する数式とリュードベリの公式}

しかし、バルマーの式が単なる「数合わせ」にとどまらなかったのは、ここからでした。

物理学者たちはこの数式を見て考えました。「もしバルマーの式にある\(2^2\)の部分を、\(1^2\)や\(3^2\)に変えたら、別の光の波長が予言できるのではないか？」

彼らは数式の仮説に従って、人間の目には見えない紫外線や赤外線の領域を必死に観測しました。すると驚くべきことに、

分母を\(1^2\)にした数式が予言する波長に、紫外線のスペクトル（ライマン系列）が発見され、

分母を\(3^2\)にした数式が予言する波長に、赤外線のスペクトル（パッシェン系列）が発見され、

さらに\(4^2\)（ブラケット系列）、\(5^2\)（プント系列）と、数式が予言する通りの場所に、次々と未知のスペクトル線が実在することが確認されたのです。

1888年、スウェーデンの物理学者ヨハネス・リュードベリは、これらバラバラに発見されたすべての系列の式を統合し、次のような「たった一つの普遍的な方程式（リュードベリの公式）」へとまとめ上げました。

\[\frac{1}{\lambda} = R \left( \frac{1}{m^2} - \frac{1}{n^2} \right)\]

（\(R\)はリュードベリ定数、\(m, n\)は整数で\(m < n\)。\(m=1\)ならライマン系列、\(m=2\)ならバルマー系列となる）

一部の可視光のデータから見つけ出したパターンが、紫外線から赤外線に至るまでのすべての水素のスペクトルを完璧に説明し、未知のデータを正確に予言する「普遍的な自然法則」へと昇華したのです。

\section{ボーアの原子模型：前期量子論の完成}

このリュードベリの公式と、プランクやアインシュタインの「飛び飛びのエネルギー（量子）」の概念を見事に結びつけ、「なぜ\(m\)や\(n\)という整数が現れるのか」という物理学的なWhyを解き明かしたのが、デンマークの物理学者ニールス・ボーアでした（1913年）。

\subsection{ラザフォード模型の致命的な欠陥}

8.1節で触れたように、当時の物理学では「中心にプラスの原子核があり、その周りをマイナスの電子が回っている（ラザフォード模型）」と考えられていました。しかし、この太陽系のようなモデルには、古典物理学的に絶対に説明できない致命的な欠陥がありました。

第4章のマクスウェル方程式によれば、電子のような電荷を持った粒がカーブして回る（加速運動をする）と、必ずそこから電磁波（光のエネルギー）が放射されます。エネルギーを放出し続ければ、電子は勢いを失い、螺旋を描いてアッという間（約100億分の1秒）に原子核に墜落して原子は消滅してしまうはずなのです。

しかし現実には、私たちを作る原子は極めて安定して存在しています。古典力学と電磁気学は、原子の内部において完全に破綻してしまったのです。

\subsection{ボーアの狂気の仮説：ツギハギの宇宙ルール}

この絶望的な矛盾に対し、ボーアはニュートンやマクスウェルの常識を一時的に「無視」するという、極めて大胆な仮説（ボーア模型）を打ち立てました。

定常状態の仮定：電子は原子核の周りを回っているが、どんな軌道でも回れるわけではない。「特定の決まった軌道（エネルギー準位）」にいる限り、電子は絶対に電磁波を放出せず、永遠に安定して回り続けることができる。

量子条件（軌道のデジタル化）：その許される特別な軌道とは、電子の運動の勢い（角運動量）が、プランク定数から作られたある単位（\(\frac{h}{2\pi}\)）の\textbf{「きっちり整数倍（\(n\)倍）」}になる軌道だけである。中途半端な軌道は存在しない。

振動数条件（量子ジャンプ）：電子が上の軌道（エネルギー\(E_n\)）から下の軌道（エネルギー\(E_m\)）へと\textbf{「瞬間移動（量子ジャンプ）」したときだけ、そのエネルギーの差額を「光の粒（フォトン）」として放出する。その光のエネルギーはアインシュタインの公式に従い、\(h\nu = E_n - E_m\)} となる。

ここでいう「軌道」は、太陽の周りを惑星が回るような実在の通り道というより、前期量子論が古典力学の言葉で無理に描いた模型です。現代の量子力学では、電子は小さな玉が決まった線の上を回るのではなく、波として広がった「軌道（オービタル）」で記述されます。ボーア模型は完全な最終理論ではありませんが、水素スペクトルの整数構造を初めて物理的に説明した、重要な足場でした。

\subsection{バルマー系列の「暗号解読」と奇跡の一致}

ボーアのこの奇想天外なモデルを使って、電子の持つエネルギー（\(E_n\)）を計算してみると、エネルギーの大きさは「軌道の番号\(n\)の2乗に反比例する（\(-\frac{1}{n^2}\)）」という結果が導き出されました。

これを先ほどの振動数条件（\(h\nu = E_n - E_m\)）に当てはめると、放出される光の波長を求める数式の中に、見事に\(\left( \frac{1}{m^2} - \frac{1}{n^2} \right)\)という形が自動的に現れました。

さらに、電子の質量や電荷、プランク定数\(h\)などの既知の数値を当てはめて計算すると、バルマーやリュードベリが「単なるデータからの数合わせ」で見つけ出していたリュードベリ定数\(R\)の値と、誤差の範囲内で完璧に一致したのです。

データ駆動で見つけ出された数式の\(m\)や\(n\)という「パズルの数字」の正体は、\textbf{「電子が飛び移る軌道の番号（量子数）」}だったのです。

数字の羅列から帰納的に導かれたリュードベリの公式が、ボーアの量子仮説によって演繹的な意味（Why）を与えられ、物理学の真理として結実した歴史的な瞬間でした。

\subsection{前期量子論の限界と、次なるパラダイムへ}

「自然界のエネルギーは連続ではなく、飛び飛び（デジタル）である」。このボーアの原子模型の成功により、量子論は物理学の主流へと一気に躍り出ました。

しかし、ボーアの理論（前期量子論）は、まだ完全なものではありませんでした。最大の謎は、\textbf{「なぜ電子は、角運動量が\(\frac{h}{2\pi}\)の整数倍になる軌道しか通ってはいけないのか？」}という点です。ボーア自身にもその根本的な理由は分かっておらず、古典力学のキャンバスに量子のルールを無理やり貼り付けた「ツギハギの理論」に過ぎなかったのです。

この「なぜ特定の軌道しか許されないのか」という究極のWhyが解き明かされるためには、第5章で学んだ「波の定在波」の概念と、物質そのものが波として振る舞うという、さらに常識外れな次なるパラダイムシフト（本格的な量子力学の幕開け）を待たねばなりませんでした。

\section{AIによる物理法則の再発見：「シンボリック回帰」の衝撃}

さて、ここで現代のAIに目を向けてみましょう。

バルマーが少数のスペクトルデータから数式を見つけ出し、それが未知のライマン系列などを予言するリュードベリ公式へと汎化していったプロセスは、まさに現代のデータサイエンスが得意とする「データからのパターン発見（帰納法）」の歴史的な大成功例です。

そして現在、AIの進化は単に「犬の画像を犬と分類する」といった予測を超え、\textbf{「実験データから、人間の物理学者が読める『物理法則の数式』そのものを自ら導き出す」}という恐るべき段階に達しています。

この技術を\textbf{「シンボリック回帰（Symbolic Regression）」}と呼びます。

\subsection{ブラックボックスを「数式」に翻訳する}

通常のディープラーニング（ニューラルネットワーク）は、何億個ものパラメータを微調整することで極めて高い予測精度を叩き出しますが、その計算過程はブラックボックス化してしまい、人間が理解できる「方程式（Why）」を出力してくれません（第12章で詳述します）。

しかしシンボリック回帰のAIはアプローチが全く異なります。このAIは、遺伝的プログラミング（生物の進化を模倣したアルゴリズム）や特殊な探索手法を用いて、変数（\(x, y, m\)など）や数学の記号（\(+\),\(-\),\(\times\),\(\div\),\(\sin\),\(\cos\), 指数など）をブロックのようにランダムに組み合わせます。そして、入力された実験データに最もよく当てはまる「数式の木構造」を、突然変異や交配を繰り返しながら探索していくのです。

出力されるのはブラックボックスな重み付けではなく、\(F=ma\)のような\textbf{「人間が読んで理解できる、記号（シンボル）で書かれた方程式」}です。

\subsection{物理学の知恵を組み込んだ「AI Feynman」}

ただランダムに記号を組み合わせるだけでは、複雑な物理現象を解くには計算量が爆発してしまいます。そこで近年、AIの探索アルゴリズムに「物理学者の知恵（公理）」を組み込むことで、劇的なブレイクスルーが起きています。

2019年、MIT（マサチューセッツ工科大学）のマックス・テグマーク教授らの研究チームは、\textbf{「AI Feynman（AIファインマン）」}と呼ばれるシンボリック回帰のシステムを発表しました。

彼らはAIに、ただ闇雲に数式を作らせるのではなく、物理学者が無意識に使っている「絶対的なルール」を教え込みました。

次元解析：「長さ（メートル）と重さ（キログラム）を足し算するような数式は、物理的にあり得ないから最初から捨てる」というルール。

対称性の利用：「宇宙のどこで実験しても（並進対称性）、単位のスケールを変えても法則は変わらないはずだ」というルール。

AI Feynmanにこれらの物理学の勘所を持たせた上で、有名なファインマン物理学の教科書に載っている「100個の複雑な物理方程式のデータ」を入力しました。するとAI Feynmanは、その100個すべての方程式（相対性理論や量子力学の式を含む）を、データのみから\textbf{再発見}してしまったのです。

もちろん、これは「AIがどんな実験データからでも自動的に真理へ到達する」という意味ではありません。AI Feynmanの問題設定では、既知の物理方程式から作られたきれいなデータを相手にしており、実際の実験データに含まれるノイズ、測定誤差、未観測の変数、装置由来の偏りはずっと厄介です。それでも、物理学の制約をうまく組み込めば、AIが単なる近似器ではなく「読める数式」を探す道具になり得ることを示した点で、この成果は大きな意味を持ちます。

\subsection{オッカムの剃刀：AIが「数式の美しさ」を評価する}

さらに、シンボリック回帰のAIは、物理学者たちが伝統的に持っている「美学」までもアルゴリズムに実装しています。それが\textbf{「オッカムの剃刀（必要なしに多くのものを定立してはならない）」}の原則です。

データに数式を当てはめようとするとき、記号を何百個も使って複雑怪奇な数式を作れば、手元のデータに誤差ゼロで完璧に一致させる（過学習する）ことは簡単です。しかし、それは「真の物理法則」ではありません。プトレマイオスの周転円のように、美しさが欠如しているからです。

シンボリック回帰のAIは、無数に生成した数式を\textbf{「データへの当てはまりの良さ（精度）」と「数式のシンプルさ（短さ）」}という2つの軸で同時に評価します。そして、最もシンプルでありながら最も高い精度を叩き出す「パレート境界（パレート最適）」に位置する数式だけを、究極の答えとして人間に提示するのです。

AIは「複雑すぎる数式は醜いから捨てる」という、アインシュタインやディラックが持っていたような「数学的な美しさへの本能」を、最適化アルゴリズムとして獲得しています。

\subsection{「AI駆動型科学」の真の幕開け}

例えば、振り子の動きや二重振り子のカオスな動きのデータをAIに入力すると、AIはニュートンの運動方程式やハミルトニアン（エネルギー保存則の数式）を自力で導き出します。さらに最近では、未知の天体の動きや、新素材の電子の動きのデータから、\textbf{「人間の物理学者がまだ誰も気づいていなかった、全く新しい物理法則の数式」}を抽出した例も報告され始めています。

バルマーやプランクが膨大なデータと格闘し、天才的な直感と「数合わせ」によって数ヶ月〜数年かけて導き出し、予測の正しさを検証していった数式を、現代のAIはわずか数時間で再発見できるようになったのです。

もちろん、AIが出してきた数式（What/How）が「なぜその形になるのか（Why）」、その背後にある物理的意味や「飛び飛びの量子条件」のような新しい宇宙のルールを解釈するのは、依然として人間の物理学者の仕事です。

物理学者が理論（演繹）から導き出した方程式の正しさをデータで検証する時代から、AIがデータ（帰納）から未知のホワイトボックスな方程式を提示し、物理学者がその意味に名前を与え、新たなパラダイムを創り出すという、新しい「AI駆動型科学（AI-driven Science）」の時代が、今まさに幕を開けようとしています。

次章では、この「前期量子論」の矛盾と限界を乗り越え、物質が「波」として振る舞うという究極の真理へと到達する「量子力学の完成（シュレディンガー方程式）」のドラマ、そしてAIにおける確率的サンプリングの共鳴へと進みます。

\section*{【第8章 章末ディスカッション課題】}

AI（シンボリック回帰）が、人類がまだ知らない全く新しい物理現象のデータから「完璧に予測が当たる数式」を導き出したとします。しかし、なぜその数式になるのか、人間の物理学者には全く理解できません。この数式を、私たちは新しい「物理法則」として教科書に載せるべきでしょうか？


\section*{参考文献（学術誌）}
\begin{enumerate}[leftmargin=2.4em]
\item Planck, M. ``Ueber das Gesetz der Energieverteilung im Normalspectrum.'' \textit{Annalen der Physik}, 4, 553--563, 1901. DOI: \url{https://doi.org/10.1002/andp.19013090310}.
\item Einstein, A. ``Uber einen die Erzeugung und Verwandlung des Lichtes betreffenden heuristischen Gesichtspunkt.'' \textit{Annalen der Physik}, 17, 132--148, 1905. DOI: \url{https://doi.org/10.1002/andp.19053220607}.
\item Bohr, N. ``On the constitution of atoms and molecules.'' \textit{The London, Edinburgh, and Dublin Philosophical Magazine and Journal of Science}, 26(151), 1--25, 1913. DOI: \url{https://doi.org/10.1080/14786441308634955}.
\item Millikan, R. A. ``A direct photoelectric determination of Planck's h.'' \textit{Physical Review}, 7(3), 355--388, 1916. DOI: \url{https://doi.org/10.1103/PhysRev.7.355}.
\end{enumerate}


